\section{Validation Gaps and Robotics Experiments}
\label{sec:discussion}
\subsection{Numerical Correctness of the Deployed Detector}
The highest-priority experiment is validation through the exact serialized
strict-DLA engine. A conventional backend that allocates linear buffers from
logical tensor shapes cannot be assumed to support DLA-native formats. The
validator must use format-aware allocation and packing, decode raw outputs with
the same class and coordinate conventions, and preserve the official COCO image
and category mappings. AP, AP$_{50}$, AP$_{75}$, and small/medium/large-object
AP should be reported for stock PyTorch, adapted PyTorch, adapted GPU FP16,
and strict DLA FP16 on the same 5,000 validation images.

Detection validation must retain low-confidence candidates and the evaluator's
multi-label and maximum-detection settings. The confidence-0.25 visual smoke
test is not a suitable AP protocol. Numerical checks should first compare raw
maps and decoded predictions on matched preprocessed inputs, then report
aggregate AP changes. This separates layout errors from the compiler's softmax
approximation and ordinary FP16 rounding.
\missing{P3}{fill strict-engine accuracy, raw-output errors, and decoder parity}.

\subsection{Ablations for the Model and Interface}
The observed compatibility result combines two changes. A controlled ablation
should use stock attention or C2DLA, crossed with decoded or raw output. Stock
attention with raw outputs isolates the attention blocker; C2DLA with decoded
outputs isolates the export boundary. These controls require fresh compiler
records rather than predictions based only on source operators. GPU accuracy
controls remain meaningful even when their strict builds fail.

A second comparison should hold the training schedule fixed while varying
shape-matched initialization, block-aware transfer, and random initialization
of the replacement block. At least three seeds would expose sensitivity to the
new mixer. A convolution-only replacement with comparable capacity would test
whether channel softmax adds accuracy worth its compiler constraints.
\missing{P5}{report build outcome, AP, parameters, and timing for these controls}.
The current single training history cannot isolate those effects or establish
that C2DLA is an optimal replacement.

\subsection{GPU Contention and Robotic Replay}
A strict learned graph does not establish isolation of the whole application.
The CPU, memory controller, power envelope, transfers, and runtime remain shared.
The first experiment should measure a fixed GPU workload alone and beside each
perception placement: stock GPU, adapted GPU, stock fallback, and adapted strict
DLA. For a co-running task with latency quantile $Q_p$, define
\begin{equation}
 S_p=\frac{Q_p(L_{\mathrm{co\text{-}running}})}
             {Q_p(L_{\mathrm{standalone}})}.
 \label{eq:slowdown}
\end{equation}
Report both median and P95 slowdown, perception latency, delivered frame rate,
and deadline misses. Use independent launches in randomized order, retaining
per-frame timestamps and clock/temperature traces. A synthetic CUDA workload
provides a repeatable stress test but does not establish a robotics benefit.

A second experiment should replay a timestamped camera sequence through a ROS~2
perception node \cite{macenski2022ros2}, alongside a specified GPU mapping or
navigation component. Keep image order, frame-release times, queue capacity, and
drop policy fixed across placements. Report acquisition-to-result latency,
queue delay, dropped frames, perception deadline misses, and co-running-task
slowdown. Compute service time alone omits the delay that a controller observes.
\missing{P6}{insert measured contention and robotic-replay results}.

\subsection{Energy at a Matched Delivered Rate}
DLA placement may change power, but longer execution and additional host work
can offset a lower device power level. Energy must therefore be integrated over
the same measured interval as completed frames:
\begin{equation}
 E_{\mathrm{frame}}=\frac{\int_{t_0}^{t_1}P_{\mathrm{board}}(t)\,dt}
 {N_{\mathrm{completed}}}.
 \label{eq:energy}
\end{equation}
Report total board energy and, separately, idle-subtracted energy using an idle
measurement under the same power mode and clock policy. Compare placements at
a common delivered rate that all can sustain, as well as under saturation;
otherwise a lower frame count can obscure the comparison. Record the rail
being sampled, its sampling interval, and alignment with runtime timestamps.
Avoid adding overlapping board and component power rails.
\missing{P7}{insert power, joules/frame, achieved rate, and repeat variability}.

\subsection{Limits of the Present Evidence}
The hardware evidence covers one Orin configuration, TensorRT 10.3, FP16, one
input resolution, and one detector scale. It establishes neither INT8 behavior
nor portability across DLA generations. Generated-input microbenchmarks exclude
content-dependent host decoding; the provisional repeated-image experiment
samples only one scene. No claim of energy savings, global GPU isolation,
hard real-time behavior, or accuracy preservation follows from strict
compilation alone. These limits define what the additional experiments must
resolve rather than changing the compatibility result already observed.
